\documentclass[11pt]{scrartcl}

\usepackage[stixtwo]{fontsetup}
\usepackage{amsmath}
\usepackage{tikz}

\begin{document}
\begin{enumerate}
	\item Bővítsük a törtünket a nevező konjugáltjával, majd végezzük el a
	      megfelelő egyszerűsítéseket:
	      \[
		      \frac{3 + 5i}{2 + 6i}
		      = \frac{(3 + 5i)(2 - 6i)}{(2 + 6i)(2 - 6i)}
		      = \frac{6 - 18i + 10i - 30i^2}{2^2 + 6^2}
		      = \frac{36 - 8i}{40}
		      = \frac{9}{10} - \frac{1}{5} i
		      \text.
	      \]
	      A számok a komplex számsíkon:
	      \begin{center}
		      \begin{tikzpicture}[scale=0.55, >=stealth]
			      \draw[->] (-1, 0) -- (4.2, 0) node[right] {$\operatorname{Re}$};
			      \draw[->] (0, -1.2) -- (0, 6.8) node[above] {$\operatorname{Im}$};
			      \fill (3, 5) circle (2pt) node[right] {$3 + 5i$};
			      \fill (2, 6) circle (2pt) node[left] {$2 + 6i$};
			      \fill (0.9, -0.2) circle (2pt) node[below right] {$\frac{9}{10} - \frac{1}{5} i$};
		      \end{tikzpicture}
	      \end{center}

	\item \textbf{Szöveges gondolatmenet}

	      A $z = 8i$ egy olyan komplex szám, amelynek valós része
	      $\operatorname{Re} \{z\} = 0$, képzetes része pedig
	      $\operatorname{Im} \{z\} = 8 > 0$.

	      Tehát a hossza $r = |z| = 8$, argumentuma pedig
	      $\arg \{z\} = 90^\circ = \pi / 2$.

	      A harmadik gyökök hossza $\sqrt[3]{8} = 2$ lesz, az egyik gyök
	      argumentuma pedig pont az eredeti szám argumentumának harmada, vagyis
	      $\varphi_1 = 90^\circ / 3 = 30^\circ = \pi / 6$. A másik két
	      gyököt az előbb meghatározott gyök $120^\circ = 2\pi/3$, illetve
	      $240^\circ = 4\pi / 3$-mal való elforgatásával kaphatjuk meg. Tehát
	      a gyökök:
	      \[
		      \sqrt[3]{8i} = \begin{cases}
			      2 \cdot e^{i \cdot \frac{\pi}{6}}
			      \text,
			      \\
			      2 \cdot e^{i \cdot \left(\frac{\pi}{6} + \frac{2\pi}{3}\right)}
			      = 2 \cdot e^{i \cdot \frac{5\pi}{6}}
			      \text,
			      \\
			      2 \cdot e^{i \cdot \left(\frac{\pi}{6} + \frac{4\pi}{3}\right)}
			      = 2 \cdot e^{i \cdot \frac{3\pi}{2}}
			      = -2i
			      \text.
		      \end{cases}
	      \]


	      \textbf{Másik megoldás}

	      $8i = 8\left(\cos \frac{\pi}{2} + i \sin \frac{\pi}{2}\right)$, így
	      \[
		      w_k
		      = \sqrt[3]{8}\left(
		      \cos\left(\frac{\pi/2 + 2k\pi}{3}\right)
		      + i \sin\left(\frac{\pi/2 + 2k\pi}{3}\right)
		      \right),
		      \quad k \in \{0; 1; 2\}
		      \text.
	      \]
	      A gyökök tehát:
	      \[
		      \begin{aligned}
			      w_0 & = 2 \left(
			      \cos \frac{\pi}{6} + i \sin \frac{\pi}{6}
			      \right)
			      \text,
			      \\
			      w_1 & = 2 \left(
			      \cos \frac{5\pi}{6} + i \sin \frac{5\pi}{6}
			      \right)
			      \text,
			      \\
			      w_2 & = 2 \left(
			      \cos \frac{3\pi}{2} + i \sin \frac{3\pi}{2}
			      \right) = -2i
			      \text.
		      \end{aligned}
	      \]
	      A gyökök egy origó középpontú, $2$ sugarú körbe írt szabályos háromszög
	      csúcsai:
	      \begin{center}
		      \begin{tikzpicture}[scale=0.8, >=stealth]
			      \draw[->] (-2.8, 0) -- (2.8, 0) node[right] {$\operatorname{Re}$};
			      \draw[->] (0, -2.8) -- (0, 2.8) node[above] {$\operatorname{Im}$};
			      \draw[dashed] (0, 0) circle (2);
			      \node[below right] at (2, 0) {$2$};
			      \draw[thin] (30:2) -- (150:2) -- (270:2) -- cycle;
			      \fill (30:2) circle (2pt) node at (42:2.5) {$w_0$};
			      \fill (150:2) circle (2pt) node at (162:2.5) {$w_1$};
			      \fill (270:2) circle (2pt) node at (282:2.5) {$w_2$};
		      \end{tikzpicture}
	      \end{center}
\end{enumerate}
\end{document}
